% ============================================================
\chapter{Introduction}
\label{ch:introduction}
% ============================================================

\section{Motivation}

Autonomous driving on a paved road is supported by a set of regularities:
lanes have conventional widths, curbs provide visible boundaries, signs encode
rules, and the road surface is designed for ordinary vehicles. An off-road
vehicle encounters none of these guarantees. A pale strip in an image may be
compacted gravel, loose soil, shallow water, or a patch of dry vegetation.
The border between a track and its surroundings can be gradual, interrupted,
or hidden by shadow. Even if the surface category were recognised perfectly,
the answer to ``can the vehicle cross it?'' would still depend on clearance,
traction, wheelbase, mass, slope, and current conditions.

Traversability is therefore an \emph{affordance}: it describes what a particular
platform can do in a scene. The same region may be accessible to a pickup,
outside a sedan's capability, and trivial for a dedicated off-road vehicle. This
vehicle dependence makes traversability more operational than ordinary
terrain classification. It also creates a demanding perception problem,
because some of the relevant physical properties are not directly observable
from one RGB frame~\cite{islam2022review,sevastopoulos2022survey}.

The practical output considered in this thesis is a dense binary mask. Every
image pixel is labelled either traversable or untraversable for a declared
vehicle-capability policy. A mask is convenient for downstream processing: a
planner can intersect it with a projected footprint, inflate blocked regions,
or combine it with geometric and temporal estimates. Dense prediction also
exposes the spatial structure of errors. A connected false-safe region directly
in front of the vehicle has a different operational meaning from a handful of
isolated pixels at the top of the image, even if both contribute the same
number of errors to a global score.

Two competing engineering requirements motivate the study. First, the mask
must be accurate enough to represent narrow corridors and irregular
boundaries. Second, it must be produced fast enough for an onboard navigation
loop. A large pretrained transformer may supply strong visual features, while
a purpose-built real-time network may offer far lower latency. The thesis
therefore studies accuracy, asymmetric error rates, reproducibility, and
hardware speed together rather than treating any one of them as a complete
measure of success.

\Needspace{0.42\textheight}
\section{Problem Definition and Research Questions}

For an RGB image
\begin{equation}
  x\in\mathbb{R}^{H\times W\times3},
\end{equation}
where $x$ is the input tensor, $H$ and $W$ are its height and width in pixels,
the final dimension contains the red, green, and blue channels, and
$\mathbb{R}$ denotes real-valued samples. The system estimates a traversable
probability for each pixel and produces
\begin{equation}
  \hat y_{ij} =
  \mathbb{1}\!\left[p(y_{ij}=1\mid x)\geq \tau\right],
  \qquad
  \hat y\in\{0,1\}^{H\times W},
  \label{eq:intro_decision}
\end{equation}
Here, $i\in\{1,\ldots,H\}$ and $j\in\{1,\ldots,W\}$ index the pixel row and
column; $y_{ij}$ is its target class; $p(y_{ij}=1\mid x)$ is the model's
conditional probability of the traversable class; $\mathbb{1}[\cdot]$ is one
when its condition is true and zero otherwise; and $\hat y_{ij}$ is the
resulting binary prediction. Class 1 is traversable, class 0 is
untraversable, and $\tau$ is the fixed decision threshold. The complete
predicted mask $\hat y$ therefore has the same $H\times W$ spatial dimensions
as the image. The main PIDNet-S development uses $\tau=0.50$ and the
CaT blue+green policy: sedan- and pickup-traversable pixels are positive,
while background and off-road-vehicle-only pixels are negative. The
architecture experiment also evaluates the more conservative blue-only policy.

This definition makes the two error directions explicit. A false positive
opens a pixel that the annotation declares blocked; it is called a
\emph{false safe}. A false negative rejects a pixel that the annotation
declares traversable; it is called a \emph{false block}. The first can expose a
planner to terrain outside the selected capability envelope. The second can
remove a feasible route or cause unnecessary stopping. The costs are
asymmetric, but they are not completely described by their pixel counts:
location, connectedness, vehicle speed, and the planner all affect the
consequence.

The investigation is organised around four research questions:
\begin{enumerate}
  \item Can PIDNet-S provide a reproducible, real-time baseline for the CaT
  binary task using only the official training partition?
  \item Which of the tested training and postprocessing choices materially
  change its accuracy, false-safe behaviour, and false-block behaviour?
  \item How does frozen-encoder transfer in ROD compare with PIDNet-S under
  matched seeds, and where do both systems lie among other compact
  segmentation architectures?
  \item How do the accuracy rankings change when inference latency, target
  hardware, execution runtime, and timing scope are included?
\end{enumerate}

\section{Objectives}

The investigation has four methodological objectives:
\begin{enumerate}
  \item define a reproducible binary task from the released CaT files,
  including explicit vehicle policies and a development split that does not
  draw from the official test partition;
  \item determine how a compact PIDNet-S implementation responds to changes in
  training objective, augmentation, class weighting, resolution, and
  threshold, while preserving the complete record of the selected model;
  \item compare that model with frozen-encoder ROD and a representative group
  of compact convolutional and transformer systems under repeatable training
  conditions;
  \item relate overlap quality and asymmetric errors to observed latency on
  CPU, H100, and TensorRT execution rather than using parameter count as a
  proxy for deployability.
\end{enumerate}

\section{Research Design}

The experiments were organised in two linked stages. The first stage treated
PIDNet-S as an engineering reference. Building that model required decisions
about file pairing, target conversion, augmentation, loss terms, validation,
error reporting, and export. Exploratory runs exposed interactions between
several of those choices; a subsequent single-factor study separated them.
The selected checkpoint was retained unchanged so that its numerical,
qualitative, and runtime measurements all refer to the same parameters.

The second stage asked whether the conclusion depended on that initial
architecture. ROD introduced a pretrained transformer whose encoder remained
fixed during CaT training. A matched PIDNet--ROD experiment measured the
difference over three seeds, after which seven additional systems supplied
context for the pairwise result. The first matrix used a common 15-epoch
budget under both label policies. Its validation histories showed that this
was useful as a fixed-budget comparison but too early to represent the models'
attainable performance. A convergence-aware follow-up therefore repeated the
principal blue+green matrix with a longer ceiling and validation-based stopping.
This order separates three questions that are often conflated: how one model
was engineered into a reproducible system, how complete model recipes compare
under equal training exposure, and how that ranking changes when each run is
allowed to progress beyond the early fixed budget.

\begin{figure}[htbp]
\centering
\resizebox{\textwidth}{!}{\begin{tikzpicture}[
  node distance=0.55cm,
  stage/.style={draw,rounded corners=2pt,minimum height=1.15cm,text width=2.55cm,
    align=center,font=\small,inner sep=5pt},
  arrow/.style={-{Latex[length=2.5mm]},thick,draw=azulUC3M}
]
\node[stage,fill=softblue] (task) {\textbf{Task definition}\\label policy, audit, split};
\node[stage,fill=softgreen,right=of task] (pid) {\textbf{Reference model}\\PIDNet-S implementation};
\node[stage,fill=softorange,right=of pid] (safe) {\textbf{Model analysis}\\ablation, asymmetric errors, reproduction};
\node[stage,fill=softpurple,below=0.9cm of safe] (rod) {\textbf{Transfer model}\\ROD with frozen encoder};
\node[stage,fill=softblue,left=of rod] (wide) {\textbf{Comparative study}\\nine systems, three seeds};
\node[stage,fill=softgreen,left=of wide] (deploy) {\textbf{Runtime study}\\CPU, H100, TensorRT};
\draw[arrow] (task) -- (pid);
\draw[arrow] (pid) -- (safe);
\draw[arrow] (safe) -- (rod);
\draw[arrow] (rod) -- (wide);
\draw[arrow] (wide) -- (deploy);
\end{tikzpicture}
}
\caption[Experimental structure.]{Experimental structure. A fixed task
definition connects model development, repeatability analysis, architectural
comparison, and hardware measurement.}
\label{fig:research_progression}
\end{figure}

Figure~\ref{fig:research_progression} summarises this progression from a fixed
task definition to model development, comparison, and hardware evaluation.

The validation-selected PIDNet-S checkpoint and the three-seed comparison are
therefore reported separately. The former is an artifact-level result backed
by exact reproduction and deployment measurements. The latter estimates
between-run variation under a common recipe. Combining them into one average
would erase the experimental conditions that give each number its meaning.

\section{Contributions}

The resulting contributions are:
\begin{itemize}
  \item an audited conversion of 1,812 underlying CaT samples to binary
  traversability, with an unchanged 544-image official test partition excluded
  from training and validation splitting and a
  deterministic 1,002/266 train--validation split;
  \item a detailed implementation of PIDNet-S for two-class prediction,
  including its tensor shapes, objective, optimisation, decision threshold,
  and controlled ablation history;
  \item a 7,623,522-parameter PIDNet-S checkpoint with 0.893941 mIoU,
  4.3176\% FSR, 7.0369\% FBR, and an exactly reproduced confusion matrix
  covering 133,693,440 pixel decisions;
  \item percentile-based and failure-ranked qualitative evidence that connects
  the aggregate asymmetric error rates to visible spatial behaviour;
  \item a CaT implementation of ROD with a frozen EfficientSAM ViT-S encoder
  and an explicit reconstruction of the published training regime;
  \item a nine-system, three-seed comparison under both blue-only and
  blue+green policies, including a $2\times2$ FPN/U-Net by
  MobileNetV2/EfficientNet-B0 design;
  \item a convergence-aware three-seed follow-up for the principal blue+green
  policy, with 27 validation-selected checkpoints, complete test records, and
  checkpoint identities;
  \item hardware measurements that distinguish forward-only CPU and H100
  timing from a larger TensorRT preprocessing-to-returned-mask scope.
\end{itemize}

\section{Scope and Responsible Interpretation}

This thesis studies monocular RGB segmentation. It does not estimate depth,
slope, soil strength, wheel slip, water depth, or a dynamically feasible
trajectory. The term \emph{safe} is retained in the metric name ``False Safe
Rate'' because it describes the binary label convention, not because the
network output certifies safe vehicle motion. A deployed system would require
sensor fusion, uncertainty handling, planner constraints, and vehicle-level
validation.

The quantitative claims are restricted to CaT, the two stated binary policies,
the declared seeds, and the saved protocols. Frames from the external ALICE
mobile-robot sequences are used only for qualitative transfer because no
ground truth is available. Hardware rates are not interchangeable: changing
the processor, runtime, numerical precision, batch size, input resolution, or
the start and end of the timer changes the meaning of FPS.

The CaT labels are treated as the reference for scoring, but not as a perfect
measurement of physical reality. Some boundaries are difficult to justify
from RGB appearance alone and may depend on physical information or subjective
vehicle-capability judgement unavailable to the model and thesis author.
Specific examples and the limits of that observation are discussed in
Chapters~\ref{ch:cat_pipeline},~\ref{ch:results}, and~\ref{ch:discussion}.

\section{Thesis Organisation}

Chapter~\ref{ch:background} introduces traversability perception, relevant
datasets, segmentation architectures, and CaT literature.
Chapter~\ref{ch:cat_pipeline} explains the CaT release, binary label policies,
split, preprocessing, and common software pipeline.
Chapter~\ref{ch:pidnet} presents the PIDNet-S implementation, exploratory
experiments, and controlled ablation.
Chapter~\ref{ch:rod_extension} introduces ROD and defines the comparative
architecture study.
Chapter~\ref{ch:protocols} defines training, evaluation, qualitative,
reproducibility, and timing protocols.
Chapters~\ref{ch:results} and~\ref{ch:hardware} report accuracy,
safety-oriented error metrics,
qualitative behaviour, and hardware speed.
Chapter~\ref{ch:discussion} analyses the trade-offs, limitations, and future
work; Chapter~\ref{ch:conclusions} summarises the evidence in relation to the
research questions.
Detailed configurations and reproduction commands are provided in the
appendices.
